% =============================================================================
% CS224R Final Report
% Hybrid Advantage Shaping with Goal-Aware Attention for Per-Turn Credit
% Assignment in LLM-Agent Reinforcement Learning
% Build: latexmk -pdf main.tex   (pdflatex + bibtex)
% =============================================================================
\documentclass{article}

\usepackage[preprint]{cs224r_2026}
\usepackage[utf8]{inputenc}
\usepackage[T1]{fontenc}
\usepackage{hyperref}
\usepackage{url}
\usepackage{booktabs}
\usepackage{amsfonts}
\usepackage{nicefrac}
\usepackage{microtype}
\usepackage{xcolor}
\usepackage{amsmath}
\usepackage{amssymb}
\setlength{\fboxsep}{2pt}
\usepackage{graphicx}
\usepackage{algorithm}
\usepackage{algpseudocode}
\usepackage{float}

\newtheorem{proposition}{Proposition}

% Avoid stretched vertical gaps around section headings (send slack to page bottom)
\raggedbottom

\title{Hybrid Advantage Shaping with Goal-Aware Attention for Per-Turn Credit Assignment in LLM-Agent Reinforcement Learning}

\author{%
  Joseph Li \\
  Stanford University \\
  \texttt{shoupei@stanford.edu} \\
  \And
  Max Rodriguez \\
  Stanford University \\
  \texttt{maxrod@stanford.edu} \\
  \And
  Samantha Leventis \\
  Stanford University \\
  \texttt{samanthaleventis@stanford.edu} \\
}

\newcommand{\Atraj}{A_{\text{traj}}}
\newcommand{\Aturn}{A_{\text{turn}}}
\newcommand{\AH}{A_{\text{H}}}
\newcommand{\flatgrpo}{flatGRPO}

\begin{document}

% ============================================================================
% PART 1: ONE-PAGE EXTENDED ABSTRACT
% ============================================================================
\begin{titlepage}
\begin{center}
\Large\textbf{Extended Abstract}\\[2pt]
\normalsize Hybrid Advantage Shaping with Goal-Aware Attention for Per-Turn Credit Assignment in LLM-Agent RL\\
\small Joseph Li, Max Rodriguez, Samantha Leventis
\end{center}
\vspace{6pt}

\paragraph{Motivation.}
LLM agents acting over many turns in sparse-reward environments such as AlfWorld and WebShop
receive only a single \emph{terminal} reward, yet success depends on credit being assigned to
the individual turns that actually mattered. Although flatGRPO is stable, it is credit-blind
because it assigns the same group-normalized trajectory advantage to every turn. Recent
step-structured variants such as GiGPO and HoGPO recognize this limitation, but still assign
credit through fixed grouping rules or heuristic step structure rather than a learned,
task-conditioned
decomposition. Classical return decomposition addresses delayed reward, but not the
language-agent setting where the same local action can be useful or harmful depending on the
goal. The missing mechanism is a GRPO-compatible way to retain the reliable trajectory
baseline while injecting learned, goal-aware per-turn credit.

\paragraph{Method.}
We close this gap with \textbf{Hybrid Advantage Shaping (HAS)}, a one-knob framework whose
per-turn advantage is a convex blend $\AH^{t} = \alpha\,\Atraj + (1-\alpha)\,\Aturn^{t}$ of
the group-normalized trajectory advantage and a per-turn signal produced by a \emph{pluggable}
decomposer. At $\alpha{=}1$ HAS recovers flatGRPO exactly, so the framework contains the
standard baseline as a special case; $\alpha{=}0.5$ is empirically optimal in our sweeps.
Our strongest decomposer is \textbf{TurnRD with goal-aware FiLM attention}: a compact
bidirectional transformer over per-turn encoder hiddens whose attention and value heads are
modulated by Feature-wise Linear Modulation (FiLM) parameters predicted from a
per-trajectory goal embedding.

\paragraph{Implementation.}
The policy is Qwen2.5-1.5B-Instruct with a rank-32 LoRA adapter (attention + MLP targets),
trained on Modal A100-80GB GPUs. Training alternates between two decoupled stages. In the
policy stage, the current LoRA policy collects $K$ rollouts for each task; trajectory
advantages are computed from group-normalized terminal rewards, while a frozen TurnRD
checkpoint predicts per-turn credits used in the HAS blend for the PPO/GRPO update. In the
credit-model stage, after a round of policy rollouts has been written to replay, the TurnRD
checkpoint is updated in a separate training job using the cumulative replay buffer with
recency decay, then loaded by the next policy-training round. This separation keeps rollout
and policy optimization lightweight while letting the credit model improve from accumulated
experience without being updated online inside every policy-gradient step.

\paragraph{Results.}
On \textbf{WebShop}, HAS with TurnRD reaches \textbf{91.3\%} success versus 80.5\% for
flatGRPO ($+10.8$pp), and beats the Progressive ($+5.3$pp) and LLM-Judge ($+9.3$pp)
decomposers.
On \textbf{AlfWorld}, TurnRD+FiLM reaches \textbf{80\%} versus ${\sim}70\%$ for flatGRPO
($+10$pp), with goal-aware FiLM contributing $+5$pp over a no-FiLM variant (75\%).

\paragraph{Discussion.}
The results validate HAS as a general credit-assignment framework: the same $\alpha$-blend
improves over flatGRPO on both AlfWorld and WebShop, while TurnRD is the strongest decomposer
among the learned and heuristic alternatives we tested. The central pattern is that
trajectory-level GRPO provides stable group-relative ranking, but needs a learned
intra-trajectory signal to separate decisive turns from incidental ones. TurnRD supplies that
signal with attention over turns, allocating credit selectively within each completed
trajectory, and goal conditioning lets the decomposer interpret actions relative to the task
instruction rather than position alone. This learned, task-conditioned credit outperforms
heuristic progress rewards and LLM-judge labels.

\paragraph{Conclusion.}
Hybrid Advantage Shaping with TurnRD offers a practical recipe for sparse-reward LLM-agent RL:
retain GRPO's stable trajectory-level objective, and add learned per-turn credit where
flatGRPO is blind. The approach is validated across two distinct agentic benchmarks, achieving
strong gains on both AlfWorld and WebShop while preserving flatGRPO as the $\alpha{=}1$
fallback. More broadly, the results support attention-based, goal-aware turn-level credit
assignment as a useful missing component for GRPO-style optimization of multi-turn LLM agents.

\end{titlepage}

% ============================================================================
% PART 2: MAIN PAPER
% ============================================================================
\maketitle

\begin{abstract}
Sparse terminal rewards make reinforcement learning for multi-turn LLM agents fundamentally a
credit-assignment problem. Success is observed only at the trajectory level, while learning
requires knowing which turns caused the outcome. GRPO-style training uses group-relative
comparisons among rollouts of the same task, but flatGRPO assigns the same
trajectory advantage to every turn and is therefore blind to intra-trajectory structure.
Recent step-structured variants add finer-grained credit heuristics, but do not learn
task-conditioned turn-level credit for language-agent trajectories. We propose \emph{Hybrid
Advantage Shaping} (HAS), a GRPO-compatible framework that preserves the stable
trajectory-level baseline while adding learned per-turn credit through a pluggable decomposer.
A single coefficient $\alpha$ interpolates between the trajectory advantage and the decomposed
turn-level signal, recovering flatGRPO exactly at $\alpha{=}1$. We instantiate the decomposer
with \emph{TurnRD}, a goal-conditioned attention model that analyzes completed trajectories
and assigns task-dependent credit to individual turns. On WebShop, HAS with TurnRD reaches
91.3\% success, improving over flatGRPO by $+10.8$pp and over Progressive and LLM-Judge
decomposers by $+5.3$pp and $+9.3$pp. On AlfWorld, TurnRD+FiLM reaches 80\% success, $+10$pp
over flatGRPO and $+5$pp over a goal-blind TurnRD variant. These results suggest that learned,
goal-aware turn-level credit is a missing ingredient for stable and sample-efficient
GRPO-style optimization of sparse-reward LLM agents.
\end{abstract}

\section{Introduction}
Large language models are increasingly deployed as \emph{agents} that act over long horizons:
issuing tool calls, navigating web pages, and manipulating embodied environments through
sequences of natural-language actions~\citep{yao2023react}. Reinforcement learning (RL) is
attractive for these agents because it optimizes task success directly rather than a proxy such
as next-token likelihood. Yet realistic agentic environments are often sparse: the environment
returns a single scalar reward only at the end of an episode, even though the episode contains
many turns. The learner must therefore solve a credit-assignment problem and determine which
turns caused success or failure, and which were incidental.

The prevailing recipe for LLM-agent RL is Group Relative Policy Optimization
(GRPO)~\citep{shao2024deepseekmath,guo2025deepseekr1}, which avoids a learned value critic by
$z$-scoring terminal rewards across a group of $K$ rollouts of the same task. This makes GRPO
simple and stable, but its standard flat form assigns the same trajectory-level advantage to
every turn of the trajectory. As a result, every turn in a successful rollout receives positive
credit, including redundant detours or actions that merely happened to precede success.

Recent work has begun to address this failure mode by adding step-level or hierarchical
structure to policy optimization, or by training process reward models that provide denser
supervision. These directions confirm that flat trajectory credit is insufficient, but they leave
two gaps for sparse-reward LLM agents: many rely on fixed grouping or externally supplied
process labels, and they do not provide a simple way to combine learned turn-level credit with
GRPO's stable trajectory-level baseline. In language-agent tasks, this distinction matters
because the same local action can be decisive or irrelevant depending on the task instruction.

We therefore argue that the missing piece is not merely denser feedback, but a safe and
task-conditioned mechanism for inserting learned per-turn credit into GRPO without discarding
the reliable group-relative trajectory signal. We make two contributions.

\textbf{(1) Hybrid Advantage Shaping (HAS)}, a framework whose per-turn advantage is a convex
blend of the trajectory advantage and a pluggable per-turn decomposer, controlled by a single
coefficient $\alpha \in [0,1]$. At $\alpha{=}1$ HAS is \emph{exactly} flatGRPO, so it
retains the baseline objective as an available fallback; intermediate $\alpha$ injects dense
per-turn credit, with $\alpha{=}0.5$ empirically optimal across both environments.

\textbf{(2) TurnRD with goal-aware FiLM attention}, the strongest realization of the
decomposer slot: a compact bidirectional transformer over per-turn encoder hiddens that
predicts per-turn credit weights and values, modulated by FiLM~\citep{perez2018film}
parameters derived from a per-trajectory goal embedding. TurnRD is trained as a side
network on a recency-decayed replay buffer with an identifiable reward-prediction objective;
in the main runs, per-turn credits are obtained by projecting value predictions onto the
constraint that their sum equals the observed terminal return.

A single shared framework delivers \textbf{91.3\%} success on WebShop ($+10.8$pp over
flatGRPO) and \textbf{80\%} on AlfWorld ($+10$pp), with goal-aware FiLM contributing $+5$pp
on AlfWorld; learned attention-based credit beats heuristic decomposers by $+5.3$ to
$+9.3$pp on WebShop. Beyond raw numbers, we contribute mechanism evidence that FiLM is genuinely
goal-aware. The remainder of the paper covers related work (\S\ref{sec:related}), the method
and Algorithm~\ref{alg:has} (\S\ref{sec:method}), setup (\S\ref{sec:setup}), results
(\S\ref{sec:results}), discussion (\S\ref{sec:discussion}), and conclusions
(\S\ref{sec:conclusion}).

\section{Related Work}
\label{sec:related}
\textbf{GRPO and critic-free LLM policy optimization.} RLHF popularized PPO~\citep{schulman2017ppo}
for aligning language models~\citep{rlhf}. GRPO~\citep{shao2024deepseekmath,guo2025deepseekr1}
removes the learned value critic by normalizing rewards within a group of sampled completions,
making it attractive for LLM fine-tuning where value estimation is costly and unstable. Recent
systems such as DAPO~\citep{yu2025dapo} and VAPO~\citep{yue2025vapo} improve GRPO-style
optimization through better sampling, clipping, and stability mechanisms. These developments
make group-relative RL practical at scale, but primarily address optimization stability: in the
flat form, a single trajectory-level advantage is still assigned to all turns. HAS is
complementary because it keeps the group-relative trajectory baseline but changes how credit is
distributed within a trajectory.

\textbf{Credit assignment in long-horizon LLM agents.} Long-horizon agent training makes
temporal credit assignment unavoidable. ArCHer~\citep{zhou2024archer} introduces hierarchical
value estimates over agent turns, while RAGEN~\citep{wang2025ragen} analyzes how GRPO behaves
as multi-turn rollouts grow longer. Recent step-structured variants such as
GiGPO~\citep{feng2025gigpo} and HoGPO~\citep{he2026hogpo} add step- or hierarchy-level
normalization to GRPO, and StepPO-style methods~\citep{wang2026steppo} align updates to
per-step signals. These methods confirm that flat trajectory credit is insufficient, but most
rely on fixed grouping, position-level structure, or externally supplied step rewards. By
contrast, HAS exposes the granularity of credit through a single mixing coefficient and allows
the per-turn signal to be learned, heuristic, or judge-based while preserving flatGRPO as the
$\alpha{=}1$ special case.

\textbf{Return decomposition and process reward models.} Classical return decomposition such
as RUDDER~\citep{arjona2019rudder} redistributes delayed rewards across timesteps, and
hindsight credit assignment~\citep{harutyunyan2019hca,andrychowicz2017her} reweights
experience by outcome. Process reward models provide denser supervision for reasoning and
agentic tasks~\citep{wu2023verify,khalifa2025prm,chen2025rmr1}, and recent
AgentPRM/HiPER-style approaches learn step-wise progress or explicit credit models for LLM
agents~\citep{xi2025agentprm,peng2026hiper}. These methods are closely related in spirit to
TurnRD, but they often require explicit process labels, external reward-model training, or a
separate supervision pipeline. TurnRD instead learns a reusable decomposer from the agent's own
replay and terminal outcomes, then plugs into a GRPO-style policy update as a post-hoc
per-turn credit signal.

\textbf{Goal-conditioned credit for language agents.} In language-agent environments such as
AlfWorld~\citep{shridhar2021alfworld} and WebShop~\citep{yao2022webshop}, the same local
action can be useful or harmful depending on the task instruction. ReAct~\citep{yao2023react}
established prompting patterns for reasoning-and-acting agents, and trajectory optimization
approaches such as ETO~\citep{song2024eto} fine-tune agents from collected experience. We
focus on how to assign credit within those trajectories. FiLM~\citep{perez2018film} provides a
lightweight mechanism for conditioning intermediate representations by feature-wise affine
modulation. TurnRD uses FiLM to condition per-turn credit on the trajectory goal, making the
decomposer task-aware rather than position-only or goal-blind. All policy updates use
parameter-efficient LoRA~\citep{hu2022lora}.

\section{Method}
\label{sec:method}

\begin{figure}[H]
  \centering
  \includegraphics[width=\linewidth]{figures/architecture.png}
\caption{\textbf{Overview of our two contributions.} \textbf{(a)} The Hybrid Advantage
  Shaping (HAS) framework: a pluggable per-turn decomposer produces $\Aturn^{t}$, blended with
  the group-normalized trajectory advantage $\Atraj$ through a single mixing coefficient
  $\alpha$. At $\alpha{=}1$ we recover flatGRPO; $\alpha{=}0.5$ is empirically optimal.
  \textbf{(b)} TurnRD with FiLM goal-conditioning, our strongest decomposer. A
  bidirectional transformer reads per-turn encoder hiddens; a per-trajectory goal embedding
  produces FiLM parameters $\gamma,\beta$ that modulate the hiddens before an $\alpha$-head
  (attention weights) and a $V$-head (per-turn values). The auxiliary reconstruction
  $\sum_t \alpha_t V_t \approx R$ makes the attention/value factorization identifiable, while
  the policy update uses a sum-preserving projection of $V_t$ for per-turn credits.}
  \label{fig:arch}
\end{figure}

\subsection{Background: Group-relative policy optimization}
For a task, the current policy $\pi_\theta$ generates a group of $K$ trajectories; each
trajectory $\tau_i$ yields a terminal reward $R_i$. GRPO forms a group-normalized trajectory
advantage by $z$-scoring,
\begin{equation}
\Atraj(\tau_i) = \frac{R_i - \mathrm{mean}_j R_j}{\mathrm{std}_j R_j + \epsilon},
\end{equation}
and updates $\pi_\theta$ with a PPO clipped surrogate plus a KL penalty to a reference policy.
In \emph{flat} GRPO this scalar is assigned to every turn (and token) of $\tau_i$; no
intra-trajectory credit is expressed.

\subsection{Hybrid Advantage Shaping}
HAS replaces the flat assignment with a per-turn advantage that blends the trajectory baseline
with a per-turn signal $\Aturn^{t}$ from a decomposer:
\begin{equation}
\AH^{t} = \alpha\,\Atraj + (1-\alpha)\,\Aturn^{t},
\qquad \alpha \in [0,1].
\label{eq:has}
\end{equation}
The blended quantity is computed \emph{per turn} and then broadcast over that turn's tokens
for the surrogate loss. Three properties make Eq.~\eqref{eq:has} attractive. First, it is a
\emph{one-knob} interpolation: $\alpha$ trades trajectory-level variance reduction against
per-turn credit density. Second, it is \emph{backward-compatible}: at $\alpha{=}1$ the
per-turn term vanishes and HAS
reduces exactly to flatGRPO, so deployments can recover the baseline objective without
changing the rest of the training stack. Third, the decomposer is \emph{pluggable}: we support
a parameter-free Progressive decomposer (env shaping signal), an LLM-Judge decomposer, and the
learned TurnRD decomposer below.

\paragraph{Why blend, and why $\alpha{=}0.5$.}
The blend is best read as a bias--variance interpolation over \emph{where} credit is placed.
The trajectory advantage $\Atraj$ is low-variance because it ranks whole trajectories within
the group and averages over much of the per-turn stochasticity. However, it is high-bias at
the turn level, assigning identical credit to every turn of a trajectory, including turns that
were irrelevant or harmful. A pure per-turn signal $\Aturn^{t}$ has the opposite profile: it
can localize credit to the turns that mattered, but it is produced by a learned decomposer and
is therefore higher-variance, especially early in training before that decomposer has
converged. Equation~\eqref{eq:has} interpolates between these regimes with a single scalar;
our $\alpha$-sweep (\S\ref{sec:discussion}) finds the midpoint $\alpha{=}0.5$ optimal on both
environments. This suggests that neither extreme suffices alone: the trajectory-only signal is
credit-blind, while an unanchored per-turn estimate is too noisy.

\begin{proposition}[flatGRPO containment]
\label{prop:safe}
At $\alpha{=}1$ the per-turn term in Eq.~\eqref{eq:has} is multiplied by zero, so
$\AH^{t}=\Atraj$ for every turn and the per-token PPO surrogate gradient is \emph{identical}
to that of flatGRPO. HAS therefore contains flatGRPO as the $\alpha{=}1$ special case and
recovers it exactly when the per-turn signal is disabled.
\end{proposition}
\noindent Proposition~\ref{prop:safe} is a containment statement, not a guarantee that every
intermediate $\alpha$ dominates the baseline. Its practical value is operational: if a
mis-specified or under-trained decomposer destabilizes learning, one can anneal
$\alpha{\to}1$ and recover the exact flatGRPO objective without modifying the trainer. We
further keep the mixture \emph{convex} ($\alpha\in[0,1]$ with weights summing to one) rather
than permitting free reweighting. Convexity preserves the scale and sign of the advantage, so
the PPO clipping range and KL trust region behave comparably to the baseline, and it keeps the
blended per-turn advantages centered consistently with the group baseline.

\subsection{TurnRD: a goal-aware attention decomposer}
TurnRD (Figure~\ref{fig:arch}b) reads the policy's per-turn last-hidden states $h_{1:T}$
($D{=}1536$), projects them to a hidden width $H{=}128$, and encodes them with a bidirectional
transformer (2 layers, 4 heads). Two heads sit on top: an $\alpha$-head that applies a softmax
over turns to produce normalized attention weights $\alpha_t$, and a $V$-head that predicts
per-turn scalar values $V_t$. The auxiliary prediction $\hat{R}=\sum_t \alpha_t V_t$ anchors
attention to return prediction rather than arbitrary saliency. For the policy update, however,
our main implementation uses a value-projection credit signal: it shifts the active-turn
values by a constant so their sum equals the terminal reward,
\begin{equation}
\hat{r}_t = \mathrm{clip}(V_t) - \frac{\sum_{s=1}^{T}\mathrm{clip}(V_s)-R}{T}.
\label{eq:vproj}
\end{equation}
The projected credits $\hat{r}_t$ define $\Aturn^{t}$ after group normalization across the $K$
rollouts at the same turn. This projection preserves the observed return exactly while
allowing negative blame and allowing one decisive turn to receive more credit than the final
reward when compensated by negative turns.

\paragraph{Why bidirectional, non-causal attention.}
The encoder is deliberately \emph{non-causal}: each turn attends to both earlier and later
turns. This is essential because the credit due to an early turn typically depends on whether
\emph{later} turns exploit it. For example, picking up an object is valuable only if a
subsequent turn uses it to satisfy the goal. A causal model, restricted to the past, cannot see the downstream
turns that realize (or waste) an early action's payoff, and would systematically under-credit
setup actions. Non-causality is legitimate here because TurnRD is an \emph{offline} post-hoc
analyzer of completed trajectories used only to shape advantages, not the acting policy; it
never leaks future information into action selection.

\paragraph{Why a softmax attention head (selectivity).}
The $\alpha$-head normalizes attention with a softmax over turns, so $\sum_t \alpha_t = 1$.
The model therefore allocates a fixed unit budget across the trajectory. This simplex
constraint forces selectivity because raising one turn's weight must lower another's, rather
than letting the model trivially flag every turn as important. In TurnRD, this attention is
used to make reward prediction identifiable and interpretable. The final per-turn shaping
signal for the main experiments comes from the value projection in Eq.~\eqref{eq:vproj}.

\paragraph{Identifiability and conservation.}
The reconstruction term $\sum_t \alpha_t V_t \approx R$ anchors an otherwise under-determined
attention/value factorization to the single observable quantity, the terminal reward, so the
attention head cannot drift into a purely cosmetic saliency map. The value-projection step
enforces exact conservation of return for the actual shaping credits: empirically the projected
per-turn credits have maximum sum-to-$R$ drift of only ${\sim}2\times10^{-6}$ across probed
trajectories, so the credit visualizations in \ref{app:extra} decompose reward rather than
merely highlighting turns.

\textbf{Goal-aware FiLM conditioning.} A per-trajectory goal embedding is projected into FiLM
parameters $\gamma,\beta \in \mathbb{R}^{H}$, which modulate each encoder hidden state before
the attention and value heads:
\begin{equation}
\mathrm{FiLM}(h_t) = \gamma \odot h_t + \beta.
\end{equation}
The FiLM projector is zero-initialized, so the conditioned model starts as the unconditioned
decomposer at round 0. As training proceeds, the model learns to use goal information to
reshape per-turn credit.

\paragraph{Why credit must be goal-aware.}
In goal-conditioned agent environments, the same local action can have different credit
depending on the task instruction. For example, opening the fridge may be useful when the goal
is to cool an object, but irrelevant when the goal is to place a book on a desk.
learn an average credit pattern over superficially similar trajectories, which can blur the
turns that matter for a specific instruction. Conditioning on the goal lets TurnRD assign
credit relative to the task being solved, rather than only relative to position or generic
progress.

FiLM provides a lightweight way to inject this goal information. A single per-trajectory affine
modulation conditions every turn representation with small parameter overhead and no dependence
on trajectory length. Because the FiLM projector is zero-initialized, the model starts as the
unconditioned decomposer and learns goal dependence only when the data supports it.
Section~\ref{sec:discussion} reports mechanism probes, and \ref{app:extra} contrasts
goal-aware and goal-blind credit on matched AlfWorld trajectories.

\textbf{Auxiliary TurnRD objective.} TurnRD is trained to predict terminal return from
per-turn representations while exposing a per-turn credit signal. Its primary objective is the
reward-reconstruction loss
\begin{equation}
\mathcal{L}_{\text{pred}}
=
\mathcal{L}_{\text{MSE}}\!\Big(\textstyle\sum_t \alpha_t V_t,\,R\Big),
\label{eq:loss}
\end{equation}
which makes the attention/value factorization identifiable from the observed return. In the
main runs we add lightweight auxiliary terms for relative ranking, progress-prior regularization,
and per-turn value supervision when environment progress signals are available. The ranking
term is a within-batch hinge loss that encourages trajectories with larger returns to receive
larger predicted returns; it complements the absolute reconstruction loss but is not the core
definition of TurnRD.

\subsection{Putting it together}
Algorithm~\ref{alg:has} summarizes the round-level training protocol. In each round, the
current LoRA policy collects $K$ trajectories per task and writes replay records containing
per-turn hidden states, goal embeddings, terminal rewards, and any available progress signals.
A separate credit-model update consumes the cumulative replay buffer with recency weighting, so
newer policy data has larger influence while earlier rounds remain available. The resulting
TurnRD checkpoint is then used as a frozen decomposer during the next policy-update stage,
where trajectory advantages are blended with TurnRD per-turn credits according to
Eq.~\eqref{eq:has}. At $\alpha{=}1$, the per-turn term vanishes and the policy objective
reduces exactly to flatGRPO, independent of the decomposer.

\begin{algorithm}[t]
\caption{Hybrid Advantage Shaping (HAS) with a replay-trained TurnRD credit model}
\label{alg:has}
\begin{algorithmic}[1]
\Require policy $\pi_\theta$ (LoRA); rounds $N$, episodes/round $M$, rollouts/task $K$;
         mix $\alpha$; recency half-life $H$
\State Replay buffer $\mathcal{B} \gets \varnothing$; initialize TurnRD checkpoint $\phi$
\For{round $r = 0,\dots,N-1$}
  \State Sample $M$ tasks and roll out $K$ trajectories each with $\pi_\theta$
  \State Append replay records to $\mathcal{B}$ \Comment{$R$, hiddens $h_{1:T}$, \texttt{goal\_emb}, round index}
  \State Update TurnRD checkpoint $\phi$ from $\mathcal{B}$ using recency-weighted minibatches
  \For{each rollout group of size $K$}
    \State $\Atraj \gets$ group-normalized $z$-score of $R$ over the $K$ rollouts
    \State $\Aturn^{t} \gets$ per-turn credit from frozen $\mathrm{TurnRD}_\phi$
    \State $\AH^{t} \gets \alpha\,\Atraj + (1-\alpha)\,\Aturn^{t}$ \Comment{broadcast over turn $t$'s tokens}
  \EndFor
  \State $\theta \gets$ PPO/GRPO clipped-surrogate update using $\AH$ with KL penalty
\EndFor
\State \Return $\pi_\theta$
\end{algorithmic}
\end{algorithm}

\section{Experimental Setup}
\label{sec:setup}
\textbf{Policy and optimization.} All methods fine-tune Qwen2.5-1.5B-Instruct with a rank-32
LoRA adapter on attention and MLP projections~\citep{hu2022lora}. We warm-start from a
supervised checkpoint trained on oracle trajectories, so the RL phase improves an already
competent instruction-following policy rather than learning the action format from scratch.
Policy updates use $K{=}8$ rollouts per task, PPO-style clipping, a KL penalty to the reference
policy, and greedy held-out evaluation ($K{=}1$) on $n{=}100$ episodes. Runs use a single
A100-80GB on Modal, ${\sim}5$--$12$h each.

\textbf{Environments.} \emph{AlfWorld}~\citep{shridhar2021alfworld} is a text embodied
household benchmark spanning navigation, object manipulation, receptacle placement, and
object-state goals; for compute control, we train on a 1{,}000-game subset of its training
split. \emph{WebShop}~\citep{yao2022webshop} is a
web-navigation shopping benchmark built from 1{,}000 products, where the agent searches,
filters, compares product attributes, and purchases an item matching a natural-language goal.
Both environments expose sparse terminal success rewards as the evaluation metric. When
available, we use environment-specific progress signals only to instantiate or supervise
per-turn decomposers, not as the final metric optimized for reporting.

\textbf{Methods compared.} We compare supervised fine-tuning (SFT); flatGRPO ($\alpha{=}1$);
and HAS ($\alpha{=}0.5$) with Progressive, LLM-Judge, and TurnRD decomposers. Progressive uses
environment-derived progress signals, LLM-Judge uses cached turn labels, and TurnRD is our
learned goal-aware decomposer. We additionally compare TurnRD with and without FiLM
goal-conditioning to isolate the contribution of goal-aware credit. Additional experiment and
implementation details are provided in \ref{app:expdetails}.

\section{Results}
\label{sec:results}

\subsection{Quantitative evaluation}
The quantitative results support two claims. First, HAS improves over trajectory-only GRPO on
both benchmarks. Second, the learned TurnRD decomposer is stronger than heuristic or judge-based
per-turn signals. We report terminal success rate throughout, using the same greedy held-out
evaluation protocol for all methods.

\textbf{WebShop.} Table~\ref{tab:webshop} and Figure~\ref{fig:webshop} compare five methods on
WebShop. SFT alone reaches 40.0\%, while flatGRPO raises success to 80.5\%, showing that
terminal-reward RL is already a strong improvement over imitation. HAS with Progressive and
LLM-Judge reaches 86.0\% and 82.0\%, respectively, both above flatGRPO but below TurnRD. The
learned TurnRD decomposer performs best at 91.3\%, improving over flatGRPO by $+10.8$pp and
over Progressive and LLM-Judge by $+5.3$pp and $+9.3$pp. This ordering indicates that the
benefit is not merely the presence of a dense per-turn signal. A replay-trained decomposer
adapted to the policy's own trajectories provides a stronger shaping signal than generic
progress heuristics or cached judge labels.

\begin{table}[H]
  \caption{WebShop terminal success rate (\%). HAS uses $\alpha{=}0.5$; flatGRPO is
  recovered at $\alpha{=}1$. TurnRD gives the strongest per-turn decomposition.}
  \label{tab:webshop}
  \centering
  \begin{tabular}{lc}
    \toprule
    Method & Success rate (\%) \\
    \midrule
    SFT (no RL)                 & 40.0 \\
    flatGRPO ($\alpha{=}1$)    & 80.5 \\
    HAS + LLM-Judge             & 82.0 \\
    HAS + Progressive           & 86.0 \\
    \textbf{HAS + TurnRD (ours)} & \textbf{91.3} \\
    \bottomrule
  \end{tabular}
\end{table}

\begin{figure}[H]
  \centering
  \includegraphics[width=0.72\linewidth]{figures/F2_webshop_5method_bar.pdf}
  \caption{WebShop success by method. HAS with TurnRD improves by $+10.8$pp over flatGRPO and
  outperforms heuristic and judge-based decomposers.}
  \label{fig:webshop}
\end{figure}

\textbf{AlfWorld.} Table~\ref{tab:alfworld} and Figure~\ref{fig:alfworld} show the same pattern
on a longer-horizon embodied benchmark. The oracle-SFT warm start reaches 47\% success, and
flatGRPO improves this to roughly 70\%, indicating that trajectory-level RL is effective but
plateaus in the high-60s. Adding TurnRD raises held-out success to 75\%, and adding FiLM
conditioning reaches 80\%. Thus, learned per-turn credit contributes about $+5$pp beyond
flatGRPO, and goal-aware conditioning adds another $+5$pp over the goal-blind TurnRD variant.
The early-round curves overlap within uncertainty, but the TurnRD+FiLM advantage emerges later
as the credit model accumulates useful replay.

The absolute success rates should not be compared directly across benchmarks. WebShop has a
shorter interaction horizon (max 15 steps) and a more structured search/filter/purchase action
space, whereas AlfWorld allows up to 40 steps and requires longer compositional action chains.
We therefore emphasize within-environment improvements over flatGRPO rather than raw
cross-benchmark success levels.

\begin{table}[H]
  \caption{AlfWorld held-out terminal success (\%). SFT is the oracle-trajectory warm start;
  flatGRPO uses trajectory-level credit only. TurnRD and FiLM each add measurable lift.}
  \label{tab:alfworld}
  \centering
  \begin{tabular}{lc}
    \toprule
    Method & Success rate (\%) \\
    \midrule
    SFT (no RL)                    & 47 \\
    flatGRPO ($\alpha{=}1$)        & ${\sim}70$ \\
    HAS + TurnRD (no FiLM)         & 75 \\
    \textbf{HAS + TurnRD + FiLM (ours)} & \textbf{80} \\
    \bottomrule
  \end{tabular}
\end{table}

\begin{figure}[H]
  \centering
  \includegraphics[width=0.82\linewidth]{figures/F4_alfworld_trajectory.pdf}
  \caption{AlfWorld round-by-round held-out success ($n{=}100$, Wilson 95\% CI bands).
  TurnRD+FiLM reaches 80\%, while flatGRPO plateaus in the high-60s.}
  \label{fig:alfworld}
\end{figure}

\subsection{Qualitative analysis}
The credit visualizations in \ref{app:extra} provide a mechanism-level view of the quantitative
improvements. On a 50-task qualitative probe slice, the goal-aware TurnRD+FiLM policy solved
44 tasks versus 39 for the goal-blind TurnRD variant; it solved every task solved by the
goal-blind policy plus five additional tasks, with no reversals. This slice is smaller than the
$n{=}100$ held-out evaluation behind the headline numbers, so we use it as qualitative evidence
rather than a second benchmark.

Across the matched examples, TurnRD concentrates positive credit on pivotal turns, such as
locating the target object, moving it to the right receptacle, or issuing the final
goal-satisfying action, while assigning small or negative credit to detours. In the case where
the goal-aware policy succeeds and the goal-blind policy fails, TurnRD+FiLM assigns high
positive credit to instruction-relevant turns, while the goal-blind variant assigns credit to
locally plausible but ultimately wrong actions. In shared-success cases, both decomposers
identify terminal goal actions, but the goal-aware model places less credit on detours. In
shared failures, negative credit concentrates on turns that misdirect the trajectory. These
examples show how learned goal-aware credit differs from flatGRPO's uniform trajectory credit
and from heuristic progress signals, which can reward local changes that do not help the
instruction.

\section{Discussion}
\label{sec:discussion}
\textbf{Why does the hybrid blend help?} The $\alpha$-sweep supports the bias--variance
interpretation of HAS. At $\alpha{=}1$, the update preserves stable group-level ranking but
remains credit-blind within each trajectory. At very small $\alpha$, the update relies too
heavily on decomposed turn-level estimates, which are noisy early in training. The strongest
setting is the hybrid regime, $\alpha{=}0.5$, where trajectory-level ranking anchors the
update while TurnRD reshapes credit within the trajectory. This explains why HAS improves most
clearly after replay has accumulated and the decomposer becomes better calibrated.

A pilot AlfWorld sweep under a smaller LoRA-rank-16 setup supports this interpretation
(Table~\ref{tab:alpha-sweep}). Although the absolute success rates are not directly comparable
to the final rank-32 runs, the relative ordering is informative: $\alpha{=}0.5$ performs best,
while both the trajectory-only setting and more extreme mixtures underperform. We therefore use
$\alpha{=}0.5$ in the main experiments.

\begin{table}[H]
  \caption{Pilot AlfWorld $\alpha$ sweep under a smaller LoRA-rank-16 setup. Absolute numbers
  are not directly comparable to the final rank-32 runs, but the sweep supports the hybrid
  setting used in the main experiments.}
  \label{tab:alpha-sweep}
  \centering
  \begin{tabular}{lc}
    \toprule
    Setting & Success rate (\%) \\
    \midrule
    flatGRPO ($\alpha{=}1.0$) & 46 \\
    HAS + TurnRD ($\alpha{=}0.25$) & 56 \\
    HAS + TurnRD ($\alpha{=}0.50$) & 62 \\
    HAS + TurnRD ($\alpha{=}0.75$) & 58 \\
    \bottomrule
  \end{tabular}
\end{table}

\textbf{What does TurnRD add beyond dense heuristics?} The decomposer comparison isolates the
source of the gain. Progressive and LLM-Judge both provide denser turn-level signals, but
neither matches TurnRD. This indicates that the important ingredient is not density alone, but
a credit model adapted to the policy's own replay distribution. Static heuristics encode
generic notions of progress, whereas TurnRD can track which turns are predictive under the
current policy and task distribution.

\textbf{What does goal conditioning add?} The AlfWorld ablation separates learned credit from
goal-conditioned learned credit. TurnRD without FiLM already improves over flatGRPO, showing
that learned per-turn decomposition is useful. Adding FiLM gives a further $+5$pp, which is
consistent with AlfWorld goals making the same local action useful or irrelevant depending on
the instruction. Figure~\ref{fig:film} reports three mechanism probes. Zeroing $\gamma/\beta$
changes the per-turn $V$-head credit by $39.7\%$, shuffling goal embeddings worsens $V$-head
MSE by $3.16\%$ (57.6\% of samples worse), and the FiLM modulation strength grows over rounds
($\gamma$ norm $0.34 \to 0.83$). These probes show that FiLM is used by the trained model, but
the modest shuffle effect also suggests that the current goal-conditioning signal is useful
rather than a complete semantic goal reasoner.

\begin{figure}[H]
  \centering
  \includegraphics[width=\linewidth]{figures/F7_film_mechanism.pdf}
  \caption{Offline FiLM mechanism probes: (a) perturbation magnitude, (b) goal-shuffle MSE,
  and (c) $\gamma/\beta$ norm growth across rounds.}
  \label{fig:film}
\end{figure}

\textbf{Failure modes and threats to validity.} The main training risk is distribution lag
between the policy and the credit model. During policy optimization, TurnRD is frozen; after
each round, its checkpoint is updated from the cumulative replay buffer with recency weighting
and then loaded for the next round. This decoupling avoids online credit-model drift within a
policy update, but the decomposer can still be imperfectly calibrated when the policy
distribution shifts across rounds. Recency weighting mitigates stale replay, but does not
eliminate this lag. Second, the goal-shuffle probe shows goal dependence, but the effect is
modest. One plausible reason is that goal-conditioned variation in raw TurnRD values is
attenuated by the downstream credit pipeline: projected per-turn credits are group-normalized,
then blended with the trajectory advantage before entering the PPO objective. Thus, FiLM may be
useful even when the measured goal-shuffle effect is small; stronger goal-discrimination losses
or adaptive normalization could make this signal more visible. Third, all reported policy
updates are LoRA-only on a 1.5B model, leaving open whether larger or fully fine-tuned policies
would require different $\alpha$ schedules. Finally, the AlfWorld curve is noisy in early
rounds. The strongest claim supported by the data is a later-stage advantage once TurnRD has
enough replay to learn useful credit, not monotonic dominance from round 0.

\section{Conclusion}
\label{sec:conclusion}
This work studies credit assignment as a central bottleneck in sparse-reward RL for multi-turn
LLM agents. We introduced Hybrid Advantage Shaping, a GRPO-compatible framework that preserves
the stable group-relative trajectory objective while adding learned per-turn credit through a
pluggable decomposer. Instantiated with TurnRD, a goal-aware attention-based decomposer, the
framework improves success on both WebShop and AlfWorld, outperforming flatGRPO as well as
heuristic and judge-based turn-level signals.

The main takeaway is that dense credit is most useful when it remains anchored to reliable
trajectory-level comparisons. flatGRPO supplies stable group-relative ranking but cannot
distinguish decisive turns from incidental ones; TurnRD supplies that missing intra-trajectory
structure by learning from replay and conditioning credit on the task goal. The consistent gains
across two qualitatively different environments suggest that learned, goal-aware turn-level
credit is a practical missing component for GRPO-style optimization of LLM agents.

More broadly, HAS provides a modular interface for studying credit granularity in agentic RL:
$\alpha{=}1$ recovers flatGRPO, while intermediate values expose how much learned per-turn
structure should influence the policy update. This makes the framework useful not only as a
strong empirical recipe, but also as a tool for analyzing when decomposed credit helps or fails
in sparse-reward language-agent training.

\section{Team Contributions}
\begin{itemize}
  \item \textbf{Joseph Li:} Researched credit assignment for LLM agents; proposed the HAS and
  TurnRD directions; designed the main experiments; implemented training infrastructure,
  RL algorithms, replay-buffer plumbing, and goal-conditioning with FiLM; and conducted more
  than 30 experiments across AlfWorld and WebShop.
  \item \textbf{Max Rodriguez:} Fine-tuned the supervised warm-start policy, explored
  alternative modeling and optimization approaches, and helped debug and improve WebShop model
  performance.
  \item \textbf{Samantha Leventis:} Ran training and evaluation scripts and trained WebShop
  models from SFT through RL variants, including TurnRD, LLM-Judge, Progressive, and flatGRPO.
\end{itemize}

\paragraph{Changes from Proposal.}
The original proposal framed the project as H-GRPO with several turn-level reward
decomposition strategies. During experimentation, we consolidated the method around the HAS
formulation and TurnRD as the strongest learned decomposer, and we added goal-aware FiLM
conditioning to TurnRD after observing that goal-blind credit assignment was insufficient for
instruction-conditioned AlfWorld tasks. We also completed the flatGRPO baseline and expanded
the evaluation to include Progressive and LLM-Judge decomposer comparisons, making the final
report focus on which source of per-turn credit actually improves sparse-reward LLM-agent RL.

\bibliographystyle{plainnat}
\bibliography{reference}

% ============================================================================
% APPENDIX
% ============================================================================
\appendix
\renewcommand{\thesection}{Appendix~\Alph{section}}

\section{Goal-aware vs.\ goal-blind credit assignment}
\label{app:extra}
To make the effect of FiLM goal-conditioning concrete, we probe the trained TurnRD decomposer
on held-out AlfWorld trajectories and visualize its per-turn credit. Each turn shows the
observation and the agent's action; the action cell is shaded by the per-turn credit (green =
positive credit, red = blame), normalized per trajectory, with the most decisive turn marked
$\bigstar$. By construction the per-turn credits sum to the trajectory's terminal reward via
the v-projection. We contrast the \emph{goal-aware} decomposer (TurnRD\,+\,FiLM) against the
\emph{goal-blind} variant (TurnRD without FiLM) on three matched cases. The headline case,
shown first, is a task where the goal-aware policy succeeds while the goal-blind one fails;
the other two cases show a shared success and a shared failure. The brief result is that
TurnRD+FiLM solves five additional tasks on this probe slice with no losses relative to the
goal-blind variant, while the examples below illustrate the mechanism: FiLM changes where
credit and blame concentrate for the same goal-conditioned task.

\input{figures/credit_sidebyside.tex}

\section{Experiment Details}
\label{app:expdetails}
\textbf{Data and evaluation.} AlfWorld experiments use a 1{,}000-game subset of the training
split for compute control and evaluate on a held-out pool with greedy decoding. WebShop
experiments use the 1{,}000-product setup described in the main text. We report terminal
success rate for both environments; dense progress signals are used only for decomposer
construction or supervision.

\textbf{SFT warm start and policy updates.} The initial policy is Qwen2.5-1.5B-Instruct with a
rank-32 LoRA adapter trained from oracle trajectories. RL updates continue from this SFT
checkpoint and keep the base model frozen, updating only LoRA parameters on attention and MLP
projection modules. Policy updates use $K{=}8$ rollouts per task, PPO-style clipping, and a KL
penalty to a reference policy.

\textbf{Replay and TurnRD updates.} Replay records store per-turn encoder hidden states,
per-trajectory goal embeddings, terminal rewards, round indices, and any available progress
signals. The replay buffer is cumulative and recency-weighted rather than physically trimmed,
so older rounds remain available but newer policy data has larger influence. TurnRD checkpoint
updates run separately from policy optimization, and the resulting checkpoint is used as a
frozen credit model during the policy stage.

\textbf{Compute and logging.} Runs execute on Modal A100-80GB workers. Per-round rollout,
TurnRD checkpoint, policy-update, and evaluation statistics are logged to manifests, including
evaluation \texttt{pct\_success}. Table~\ref{tab:hparams} summarizes the shared
hyperparameters.

\begin{table}[H]
  \caption{Shared hyperparameters.}
  \label{tab:hparams}
  \centering
  \begin{tabular}{ll}
    \toprule
    Setting & Value \\
    \midrule
    Base model            & Qwen2.5-1.5B-Instruct \\
    LoRA rank / targets   & 32 / attention + MLP \\
    Rounds                & 10 \\
    Episodes per round    & 80 \\
    Rollouts per task $K$ & 8 \\
    Mixing coefficient $\alpha$ & 0.5 (HAS), 1.0 (flatGRPO) \\
    TurnRD encoder        & 2 layers, 4 heads, hidden $H{=}128$ \\
    Recency half-life $H$ & 4 rounds \\
    Eval                  & greedy, $n{=}100$ held-out / round \\
    Hardware              & A100-80GB (Modal) \\
    \bottomrule
  \end{tabular}
\end{table}

\end{document}
